\chapter{Meaning Games: Wittgenstein, Nash, and the Strategic Foundation of Sense}\label{ch:meaning-games}

\epigraph{For a \emph{large} class of cases---though not for all---in
which we employ the word ``meaning'' it can be defined thus: the meaning
of a word is its use in the language.}{Wittgenstein, \emph{Philosophical
Investigations} \S43~\cite{Wittgenstein1953}}

The slogan that meaning is use is one of the most quoted lines in
twentieth-century philosophy and one of the least
\emph{specified}. Anglophone commentators have read it as a banner for
deflationism, behaviourism, anti-realism, semantic externalism, social
pragmatism, and several incompatible verificationisms. The reason is
structural: ``use'' is a verb without a fixed type. Use \emph{by whom},
on \emph{what occasions}, with \emph{what payoffs}, against \emph{whose
choices}, in \emph{what space of alternatives}? Without these slots
filled, the slogan is a gesture; with them filled, it is the definition
of a non-cooperative game.

The thesis of this chapter is that those slots are filled by Idea
Theory. The carrier $\Ideas$ supplies the alternatives; the composition
$\compose$ supplies the rule by which two uses combine into a third; the
resonance $\rho$ supplies the payoff differential by which one use is
preferred to another; the agents are speakers; the equilibria are
\emph{conventions} in the sense of \cite{Lewis1969} and the focal
points of \cite{Schelling1960}. Once these identifications are in place,
the late Wittgenstein turns out to be making a precise mathematical
claim, the rule-following paradox of \cite{Kripke1982} turns out to be a
finite-underdetermination lemma, the impossibility of a private language
turns out to be the triviality of a one-player coordination game, and
Brandom's inferentialism~\cite{Brandom1994} turns out to be a Pareto
refinement on the same equilibrium concept. The chapter develops these
identifications end-to-end, drawing on two layers of machine-checked
formalization. The first is a general game-theoretic layer (verified in
\texttt{IdeaTheory.Extensions.GameTheory}, nineteen theorems) which
establishes existence and structural properties of equilibria for games
whose strategy spaces and payoff structures are induced by an idea
algebra. The second is a Wittgensteinian layer
(verified in \texttt{IdeaTheory.Extensions.MeaningGames},
fourteen-plus theorems) which specializes the general layer to
coordination on shared meanings under $\compose$. The final two
sections relate the resulting \emph{meaning games} to the macroscopic
diffusion processes of \texttt{IdeaTheory.Extensions.Diffusion} and the
recombinant growth dynamics of
\texttt{IdeaTheory.Extensions.RecombinantGrowth}, which are themselves
machine-checked.

%======================================================================
\section{Two pictures of meaning}\label{sec:two-pictures}
%======================================================================

The Tractatus~\cite{Wittgenstein1922} offered a single, austere account
of how a sentence comes to mean anything at all. A proposition is a
\emph{picture} of a possible state of affairs. The picturing relation is
itself an isomorphism: the logical multiplicity of the sign matches the
logical multiplicity of the situation. Meaning is grounded
\emph{externally}, in the rigid mapping between a sign-structure and a
world-structure. The arrows of the diagram run from signs into the
world, and the diagram is fixed once the lexicon and the syntax are
fixed.

The \emph{Investigations}~\cite{Wittgenstein1953} pulled this picture
apart from inside. Three observations did most of the work. First, the
same sign-structure supports indefinitely many uses---ostensive
definition, command, description, joke, prayer, theorem, insult---no
one of which has structural priority. Second, the meaning of a word is
not given by an inner mental act but by participation in a public
practice. Third, the practice is itself open-ended: the rules for
correctness do not determine themselves, but are determined by the
sustained agreement of the speakers (\S\S198--202).

These two pictures look incompatible because they appear to disagree
about the \emph{location} of meaning. Tractarian meaning is a relation
between a sign and a world; Investigations meaning is a relation between
a sign and a community. We propose that the disagreement is illusory once
the community is allowed to occupy the place of the world: the
constraint that fixes a sign's content is, in either case, that the
sign be \emph{coordinable}. In the Tractatus the sign coordinates with a
state of affairs; in the Investigations it coordinates with the uses of
other speakers. In both cases the formal object that fixes meaning is
a fixed point of a relation. Idea Theory makes this fixed-point
intuition mathematical: the relation is composition, the fixed points
are the equilibria of a coordination game, and the world-side and
community-side readings differ only in whether the other player is the
world or the community.

The strategic re-reading of the slogan also resolves a methodological
puzzle. If meaning is use, why is it not arbitrary? Why does
$\textit{red}$ mean red and not blue, given that we could have used
either? The answer cannot be that $\textit{red}$ pictures redness, on
pain of regress. The answer Wittgenstein in fact gives---that we form
a community in which calling certain things \emph{red} is what we
do---is, on its surface, sociological. But once \emph{what we do} is
read as a strategy profile in a game with payoffs that reward shared
composition, the answer becomes mathematical: $\textit{red}$ means red
because the profile in which speakers use the sign for that colour is a
Nash equilibrium, and the alternative profiles are not. The strategic
formalization, in short, salvages the determinacy of the Tractarian
account without its metaphysics; it salvages the openness of the
Investigations account without its mysticism. We shall see in
\S\ref{sec:meaning-game-defined} that this is not a metaphor.

There is a subtler reading on which the two pictures differ not in
\emph{location} but in \emph{cardinality of the coordinating party}.
The Tractatus picture pairs a sign with a single, unique state of
affairs; the picturing relation is a function. The Investigations
picture pairs a sign with the joint use of a community; the
coordinating relation is a many-to-many correspondence whose stable
points are the equilibria of a game. The Tractatus is the limiting
case in which the community has cardinality one and the unique world
plays the role of the second coordinator. We treat this not as a
charitable reconciliation but as a fact about the formalism: the
function-versus-correspondence distinction is recovered by varying
$|N|$ in \cref{def:meaning-game}, and the Tractarian picture survives
as the boundary case $|N|=1$ together with a degenerate $\rho$
restricted to the sign--world pairing. The collapse of meaning under
$|N|=1$ that we prove in \S\ref{sec:private} is then the formal
explanation of why the late Wittgenstein had to abandon the early
picture: the limiting case is degenerate, and the non-degenerate
account requires $|N|\geq 2$.

\formalstatus The conceptual identification of meaning with equilibrium
strategy is captured by the type signature of a meaning game in
\texttt{IdeaTheory.Extensions.MeaningGames}; the specific Tractatus
versus Investigations contrast is foundational background and is not
itself a formalized theorem.

%======================================================================
\section{Meaning games defined}\label{sec:meaning-game-defined}
%======================================================================

\begin{definition}[Meaning game]\label{def:meaning-game}
Let $(\Ideas,\compose,\unit,\rho,\sigma,\nu,\eta,\omega,\deg)$ be an
idea algebra. A \emph{meaning game} on the algebra is a tuple
$\Gamma=(N,\Ideas,M,(u_i)_{i\in N})$ where:
\begin{enumerate}[leftmargin=2em]
  \item $N$ is a non-empty finite set of \emph{speakers};
  \item $M\subseteq\Ideas$ is a non-empty subset whose elements we call
    \emph{intelligible meanings};
  \item each $u_i:\Ideas^N\to\RR$ is a payoff function on the strategy
    profile space whose elements are \emph{uses} $s=(s_i)_{i\in N}$,
    $s_i\in\Ideas$;
  \item the canonical payoff is the \emph{coordination payoff}
    \[
      u_i(s) = \rho\!\left(s_i,\ \bigdiamond_{j\in N} s_j\right)
              \cdot \mathbf{1}_{M}\!\!\left(\bigdiamond_{j\in N} s_j\right),
    \]
    where $\bigdiamond_{j\in N} s_j$ denotes the iterated $\compose$
    aggregate of the profile (taken in any fixed order; differences
    between orders are absorbed by the cocycle $\omega$ and the
    aggregate payoff is well-defined modulo $\ker\rho$).
\end{enumerate}
\end{definition}

A pure-strategy \emph{Nash equilibrium} of $\Gamma$ is a profile $s^\ast$
such that no speaker $i$ can strictly improve $u_i$ by unilaterally
substituting some $t_i\in\Ideas$ for $s^\ast_i$. A \emph{convention} of
$\Gamma$ is a Nash equilibrium whose aggregate
$\bigdiamond_{j\in N} s^\ast_j$ lies in $M$.

The definition encodes three things at once. First, that meaning is
typed: a use is an element of $\Ideas$ and only an aggregate that lands
in $M$ counts as having said anything intelligible. Second, that meaning
is \emph{joint}: payoff depends on the composite of all speakers'
contributions, not on any single speaker's act in isolation. Third,
that meaning is \emph{competitive against alternatives}: the equilibrium
condition rules out unilateral improvements, which is exactly the
mechanism by which conventional choices among logically equivalent
alternatives become stable.

The choice of $\rho$ as the payoff kernel is principled. Theorem~2 of
the foundations~\cite{Nash1950} establishes $\rho$ as the unique
bilinear extension of the resonance pairing on $\Ideas$, so the payoff
is automatically linear in mixed strategies, the game has a well-defined
mixed extension, and the standard machinery for existence of equilibria
applies. The indicator $\mathbf{1}_M$ is the formal trace of
Wittgenstein's observation that not every utterance is a move in the
language; the indicator is exactly what disqualifies a profile that
aggregates to nonsense.

\begin{theorem}[Existence of conventions]\label{thm:existence-conv}
Let $\Gamma$ be a meaning game with $|N|<\infty$ and $M\neq\emptyset$
finite. Then $\Gamma$ admits at least one (mixed-strategy) Nash
equilibrium whose aggregate is supported in $M$.
\end{theorem}

\begin{proof}[Proof sketch]
Restrict each speaker's strategy set to a finite generating set
$S_i\subseteq\Ideas$ that contains, for each $m\in M$, at least one
$\compose$-decomposition into $|N|$ factors. Theorem~1 of the
foundations guarantees that such decompositions exist (take $|N|-1$
copies of $\unit$ and one copy of $m$). The resulting finite normal-form
game has, by the standard finite-strategy existence result, a Nash
equilibrium; the indicator term in $u_i$ ensures that any equilibrium
with strictly positive payoff for some speaker has aggregate in $M$,
and the trivial convention $(s_i)\equiv\unit$ shows that an aggregate-in-$M$
equilibrium always exists when $\unit\in M$. The general case follows by
relabelling.\end{proof}

A few remarks on the definition are in order. First, the iterated
$\compose$ aggregate is not in general associative on the nose---only
associative up to the cocycle $\omega$. The payoff specification is
nevertheless well-defined because $\rho$ vanishes on coboundaries
(Theorem~5), so re-association of the aggregate changes the payoff by
an exact differential which is invisible to the equilibrium condition.
Second, the order of speakers in the aggregate is a convention; if
$\compose$ is non-commutative, different orderings yield different
aggregates, but the equilibria of a meaning game played in any fixed
order are in canonical bijection with those played in any other fixed
order, via conjugation by the appropriate permutation
representation (Theorem~6). Third, the indicator $\mathbf{1}_M$ is the
formal trace of \emph{intelligibility}: an utterance whose aggregate
fails to land in $M$ produces no payoff for any speaker, modelling the
phenomenon that nonsense produces neither agreement nor disagreement.
Fourth, the use of $\rho$ rather than a generic utility functional is
not arbitrary: it is the unique bilinear form on $\Ideas$ guaranteed
by Theorem~2, and any other bilinear payoff would be a linear
combination of $\rho$ and its curvature $\kappa$ (Theorem~5), the
latter contributing only a redistribution among speakers and not
affecting the equilibrium structure on the symmetric component.

\formalstatus \cref{def:meaning-game} and \cref{thm:existence-conv} are
verified in \texttt{IdeaTheory.Extensions.MeaningGames}; the underlying
existence machinery for Nash equilibria of finite normal-form games
induced by an idea algebra is verified in
\texttt{IdeaTheory.Extensions.GameTheory}.

%======================================================================
\section{Effective meaning as aggregate}\label{sec:effective}
%======================================================================

A convention is a profile, and a profile is a tuple of
contributions. What is publicly available---what counts as
\emph{the meaning} that the community has settled on---is not any
single contribution but the aggregate of them under composition. We
call this the \emph{effective meaning} of the convention. The following
lemma makes the identification precise.

\begin{lemma}[Effective Meaning Lemma]\label{lem:effective}
Let $s^\ast$ be a convention of $\Gamma$ with aggregate
$m^\ast\defeq\bigdiamond_{j\in N} s^\ast_j$. Then $m^\ast\in M$, and for
any speaker $i$, the marginal contribution
$\rho\!\left(s^\ast_i,\ m^\ast\right)$ is non-negative and is
maximized over $\Ideas$ by $s^\ast_i$ itself, holding fixed
$(s^\ast_j)_{j\neq i}$.
\end{lemma}

\begin{proof}
Membership of $m^\ast$ in $M$ is the convention condition. Fix $i$ and
write $a\defeq \bigdiamond_{j<i} s^\ast_j$, $b\defeq \bigdiamond_{j>i}
s^\ast_j$, so that the aggregate as a function of speaker $i$'s
strategy $t$ is $a\compose t\compose b$. Speaker $i$'s payoff is
\[
  u_i(t,s^\ast_{-i}) = \rho(t,\ a\compose t\compose b)\cdot
  \mathbf{1}_M(a\compose t\compose b).
\]
By the equilibrium condition, $u_i(s^\ast_i,s^\ast_{-i})\geq u_i(t,s^\ast_{-i})$
for all $t$. Taking $t=\unit$ gives $u_i\geq\rho(\unit,a\compose b)\cdot
\mathbf{1}_M(a\compose b)$; if $\unit$ enters $M$ this lower-bounds the
optimal payoff at $0$, and since $\rho$ extends bilinearly with $\rho(\unit,\cdot)$
fixed by the cocycle normalization $\omega(\unit,\cdot)=0$ (Theorem~4),
the marginal $\rho(s^\ast_i,m^\ast)$ is the unique attainable maximum.
\end{proof}

The lemma is innocuous-looking but does crucial work for the rest of
the chapter. It says that the publicly available object---the thing the
community is talking about---is not the lexical entry of any individual
but the composite to which all individuals' uses contribute. This is the
formal correlate of two of Wittgenstein's most quoted observations:
that one cannot unilaterally fix what a word means
(\S\S243--315) and that the word's meaning is what survives the
back-and-forth of public correction. The marginal-contribution clause
is the formal correlate of accountability: each speaker's use is
optimized against the aggregate, and a speaker who deviates pays in
$\rho$.

A second consequence of the lemma is that effective meaning is
\emph{shareable across the equilibrium}: $m^\ast$ depends only on the
profile, not on which speaker's perspective we take. This is the
property that distinguishes convention from idiolect. An idiolect is a
strategy that maximizes a single speaker's payoff under arbitrary
guesses about the rest of the profile; a convention is a profile in
which the optimization has converged for all speakers simultaneously.
The asymmetry between the two is exactly the asymmetry between
unilateral best-response and joint equilibrium. Sections
\S\ref{sec:private} and \S\ref{sec:salience} return to this point.

A third consequence concerns \emph{accountability under correction}.
When a speaker $i$ is corrected by a co-speaker---``no, that's not
what the word means''---the formal content of the correction is
information that $s^\ast_i$ would receive lower payoff than some
alternative against the prevailing aggregate. The corrector communicates
the marginal: the gradient of $u_i$ in strategy space at the current
profile points away from $s^\ast_i$. Convergence to the equilibrium is
the speaker's gradient ascent on $u_i$, holding the aggregate fixed,
followed by recomputation of the aggregate after each speaker has
adjusted. The Effective Meaning Lemma is the fixed-point statement of
this iterative correction process: at convergence, no further correction
is informative because every speaker is already at her best response.

\formalstatus The Effective Meaning Lemma is verified in
\texttt{IdeaTheory.Extensions.MeaningGames}; the underlying bilinearity
and cocycle-normalization properties of $\rho$ are Theorems~2 and 4 of
the foundations.

%======================================================================
\section{Family resemblance}\label{sec:family-resemblance}
%======================================================================

\textbf{Original claim.} \emph{PI} \S\S65--71~\cite{Wittgenstein1953}.
There is no single feature shared by all uses of \emph{game}; rather,
games form a network of overlapping similarities the way the members of
a family share various features pairwise without any feature being
shared by all. Wittgenstein's negative thesis is that the search for a
common essence is a misapprehension; his positive thesis is that the
network of overlaps is itself enough to support stable use.

\textbf{Formal restatement.} A subset $F\subseteq\Ideas$ exhibits
\emph{family resemblance} when, for every $a,b\in F$, there exists
$c\in\Ideas$ with $c=\sigma(a,c)=\sigma(c,b)\neq\unit$, i.e.\ a
non-trivial common $\sigma$-idempotent. Equivalently, the binary
relation $a\!\smile\!b\Leftrightarrow\exists c\neq\unit:c=\sigma(a,c)=\sigma(c,b)$
on $F$ is connected but not necessarily transitive. Note that we do
\emph{not} require a single $c$ to work for the whole family.

\begin{theorem}[Family-resemblance conventions]\label{thm:family-resemblance}
Let $F\subseteq\Ideas$ exhibit family resemblance and let $M$ be the
$\sigma$-closure of $F$. Then the meaning game $\Gamma=(N,\Ideas,M,u)$
admits a non-trivial convention; moreover, distinct equilibria of
$\Gamma$ may have aggregates in distinct connected components of the
$\smile$-graph on $F$, so the convention concept itself is plural.
\end{theorem}

\begin{proof}
Existence is by \cref{thm:existence-conv} applied to the non-empty
$M$. To rule out the trivial convention as the only one, fix any
$a,b\in F$ with $a\smile b$ via $c\in\Ideas$ and set $s^\ast_i=c$ for
all $i$. Then $\bigdiamond_i s^\ast_i$ lies in the $\sigma$-orbit of
$c$, hence in $M$, and unilateral substitution of $\unit$ strictly
decreases payoff because $\rho(\unit,c)=0$ by Theorem~4 normalization
while $\rho(c,c)>0$ by non-degeneracy on the support of $F$. Plurality
follows from connectedness without transitivity of $\smile$: two
components yield two non-equivalent convention aggregates.
\end{proof}

The theorem makes Wittgenstein's negative thesis a feature of the
meaning game rather than a bug. There is no single $\sigma$-invariant
common to the family; nevertheless, every pairwise overlap is enough to
support a stable equilibrium. The plurality clause is what licenses the
late Wittgensteinian observation that one and the same family-name
covers genuinely different equilibria: \emph{board game}, \emph{ball
game}, \emph{language game} can all be conventions on the same name
without any of them being a privileged one. This is a version of the
quotient construction made precise: $F$ projects onto the equivalence
classes of the $\smile$-graph, and each class supports its own
convention.

The contrast with classical theories of concept structure is sharp.
Classical definitionism (the necessary-and-sufficient-conditions theory
of \cite{FodorPylyshyn1988}) requires a single $\sigma$-invariant
shared by all instances. Prototype theory (\cite{Rosch1973,
RoschMervis1975}) replaces the invariant by a graded similarity to a
typical member, which corresponds in the formalism to ordering $F$ by
$\rho$-distance from a focal $c\in F$. Family resemblance, on our
reading, drops both: it asks neither for an invariant nor for a focal
prototype, but only for a connected pattern of pairwise overlaps. Idea
Theory makes all three positions formally available and consigns the
debate among them to the structural question \emph{which condition
holds in the algebra realizing a given conceptual domain}.

A concrete payoff of the family-resemblance theorem is that it
explains why bilingual or polysemous lexical items are stable. The
English word \emph{run} covers the running of an athlete, of a
machine, of a candidate, of a colour, of a stocking, of a melody.
Classical theory predicts collapse to a single sense or fragmentation
into homonyms; prototype theory predicts a single graded
generalization from a focal sense. Family-resemblance theory predicts
exactly what is observed: a connected but non-transitive network of
overlapping pairwise senses, each pair sharing a $\sigma$-idempotent
without there being a global one. The prediction extends to the
historical observation that polysemous words are diachronically
stable: they support distinct conventions in distinct contexts, and
the conventions stabilize independently because the equilibrium set
of the meaning game is itself plural.

\formalstatus \cref{thm:family-resemblance} is verified in
\texttt{IdeaTheory.Extensions.MeaningGames}; the underlying
$\sigma$-decomposition is Theorem~3 of the foundations.

%======================================================================
\section{The private-language argument formalized}\label{sec:private}
%======================================================================

The most controversial argument of the late
Wittgenstein---\emph{PI} \S\S243--315~\cite{Wittgenstein1953}---is that a
language whose words refer to private sensations and whose correctness
is not subject to public check is impossible. The argument has been
read deflationarily (the diarist's $S$ has no use, hence no meaning),
verificationally (no public check, no truth condition), and
behaviouristically (no behavioural criterion, no concept). What it has
not been read as---and what it is, on our analysis---is a triviality
result for one-player coordination games.

\begin{theorem}[Triviality of private convention]\label{thm:private}
Let $\Gamma=(\{1\},\Ideas,M,u_1)$ be a meaning game with a single
speaker. Then every convention of $\Gamma$ has aggregate $\unit$, and
the only Nash equilibrium with aggregate in $M\setminus\{\unit\}$ is
the all-$\unit$ profile if $\unit\in M$.
\end{theorem}

\begin{proof}
With $|N|=1$ the aggregate of any profile is the single strategy
$s_1$. The payoff is $u_1(s_1)=\rho(s_1,s_1)\cdot\mathbf{1}_M(s_1)$,
which is best-responded by maximizing $\rho(s_1,s_1)$ subject to
$s_1\in M$. The Nash condition reduces to: $s_1$ is a global maximum of
$\rho(\cdot,\cdot)$ on $M$ along the diagonal. But this maximum is
attained only by $\unit$, because the bilinearity of $\rho$ together
with the normalization $\omega(\unit,\cdot)=0$ (Theorem~4) makes
$\rho(\unit,\unit)$ the unique fixed point of the diagonal map; any
other diagonal value involves the antisymmetric curvature $\kappa$
(Theorem~5), and is therefore not a maximum. Hence the only Nash
equilibrium has aggregate $\unit$.
\end{proof}

The theorem is the Wittgenstein--Nash version of the private-language
argument. The mechanism is identical to Wittgenstein's: a single
speaker has no second party against whom to coordinate, so there is no
non-trivial alternative to coordinate \emph{on}, so the only stable
``meaning'' is the empty one. The diarist's $S$ collapses to $\unit$
because the diarist is the only player. The verificationist reading
misidentifies the absent ingredient as a public check; the strategic
reading correctly identifies it as a second player. The behaviourist
reading misidentifies the absent ingredient as observable behaviour; the
strategic reading correctly identifies it as a payoff structure with
non-trivial best responses, which a one-player game cannot have.

A further consequence is that the argument is sensitive to the
\emph{cardinality} of $N$, not to the metaphysics of mind. Two-player
meaning games already support non-trivial conventions
(\cref{thm:existence-conv}); the argument's force is therefore
combinatorial, not phenomenological. Critics who object that
Wittgenstein has misdescribed the introspective situation of the
diarist (\cite{Kripke1982} surveys several) have not engaged with the
strategic content of the argument, which is independent of any claim
about what the diarist can or cannot notice. The diarist's notebook
fails not because the entries are wrong but because there is no game
in which they could be moves.

\formalstatus \cref{thm:private} is verified in
\texttt{IdeaTheory.Extensions.MeaningGames}; it depends on the cocycle
normalization (Theorem~4) and on the symmetric-versus-curvature
decomposition (Theorem~5).

%======================================================================
\section{Salience and Lewis convention}\label{sec:salience}
%======================================================================

\cite{Lewis1969} pointed out that coordination games standardly possess
many Nash equilibria, and that the question \emph{which equilibrium
gets played} cannot be answered from inside the equilibrium concept.
Lewis's answer was \emph{salience}: an equilibrium is salient when it
has features that draw the agents' attention preferentially. He took
these features to be public common knowledge of conspicuous attributes
(precedent, simplicity, naturalness). On our reading the salient
features are the $\sigma$-idempotents inside $M$.

\begin{definition}[Salient use]\label{def:salient}
An element $c\in\Ideas$ is \emph{salient relative to $M$} when
$c\in M$ and $c=\sigma(c,c)$, i.e.\ $c$ is a $\sigma$-idempotent
intelligible meaning. The set of salient meanings is denoted
$\mathrm{Sal}(M)$.
\end{definition}

\begin{theorem}[Strict-Nash refinement at salient idempotents]\label{thm:salient}
Let $\Gamma$ be a meaning game and let $c\in\mathrm{Sal}(M)$. The
profile $s^\ast_i\equiv c$ is a strict Nash equilibrium of $\Gamma$:
every unilateral deviation $t\neq c$ strictly decreases the deviator's
payoff. Conversely, every strict Nash equilibrium of $\Gamma$ whose
aggregate is in $\mathrm{Sal}(M)$ has profile constant on
$\mathrm{Sal}(M)$.
\end{theorem}

\begin{proof}
Since $c=\sigma(c,c)$, the aggregate of the constant profile is $c$
itself (idempotence absorbs the cocycle correction by Theorem~4). The
deviator $i$, choosing $t\neq c$, changes the aggregate to
$\sigma(c,t)\compose c\compose\cdots$; by Theorem~3 the
$\sigma$-component of this is no longer $c$, and by Theorem~5 the
$\nu$ and $\eta$ components contribute non-zero curvature, so
$\rho(t,\text{new aggregate})<\rho(c,c)$. Strictness follows from
non-degeneracy of $\rho$ on the support. The converse uses uniqueness
of $\sigma$ modulo $\ker\rho$ (Theorem~3) to recover constancy.
\end{proof}

The theorem upgrades Lewis's empirical observation about salience to a
formal refinement of the equilibrium concept: among the many Nash
equilibria of a meaning game, those whose aggregate is a
$\sigma$-idempotent are \emph{strict}, hence robust to small
perturbations of strategy. The selection problem is partially
solved---inside the formalism, salient equilibria are mathematically
distinguished. Empirically, the question whether a given convention is
salient becomes the question whether its aggregate is a
$\sigma$-idempotent in the algebra realizing the language, which is in
principle measurable.

The reading also explains why salience-based conventions are stable
under perturbation in a way that mere Nash conventions are not. A
non-salient equilibrium tolerates indifferent deviations: speakers who
switch to a behaviourally equivalent alternative pay nothing, and the
equilibrium can drift. A salient equilibrium punishes every deviation
in $\rho$, so drift is suppressed. This formalizes the intuition that
``natural'' conventions---the ones that pick themselves out
spontaneously---are also the ones that resist being changed.
\cite{Schelling1960} drew the same conclusion empirically; we return to
his focal points in \S\ref{sec:focal}.

A subtle point about strict-versus-non-strict equilibria deserves
emphasis. \cref{thm:salient} shows that salient idempotents support
\emph{strict} equilibria. Not every salient meaning is a strict
equilibrium of every meaning game---only those whose surrounding
strategy space contains the requisite variation to make the deviation
loss strict. But every strict equilibrium of every meaning game has
its aggregate at a salient idempotent. The biconditional that this
suggests---salient idempotents in $M$ are in bijection with strict
conventions of $\Gamma$---holds under the additional assumption that
the support of strategy space is dense in $\Ideas$, which is the
generic condition. Lewis's empirical observation that the salient
equilibria are also the ones that survive perturbation is recovered
exactly in this generic regime.

\formalstatus \cref{def:salient} and \cref{thm:salient} are verified
in \texttt{IdeaTheory.Extensions.MeaningGames}; idempotence and
strictness rely on Theorems~3, 4, and 5 of the foundations.

%======================================================================
\section{Rule-following and Kripkenstein}\label{sec:rule-following}
%======================================================================

\cite{Kripke1982} extracted from \emph{PI} \S\S201--202 a sceptical
paradox: any finite stretch of past usage of a rule is compatible with
infinitely many incompatible continuations. Kripke's notorious
example is the function $\textit{quus}$, which agrees with $\textit{plus}$
on all argument pairs the speaker has ever computed and diverges
elsewhere. The paradox is that the speaker's past behaviour does not
fix which of $\textit{plus}$ and $\textit{quus}$ she has been
following, so the very notion of \emph{following a rule} threatens
incoherence.

\begin{lemma}[Finite underdetermination]\label{lem:underdet}
Let $H=(s^{(1)},\ldots,s^{(T)})$ be a finite history of strategy
profiles in a meaning game $\Gamma$. Then there exist infinitely many
distinct best-response correspondences
$\beta:\Ideas^N\to\mathcal{P}(\Ideas)$ each of which is consistent
with $H$ in the sense that $s^{(t)}_i\in\beta(s^{(t)}_{-i})$ for all
$t,i$, and which differ on profiles not occurring in $H$.
\end{lemma}

\begin{proof}
Define $\beta_0$ as the literal best-response correspondence of
$\Gamma$ on the support of $H$, extended trivially. For each
$\bar s\in\Ideas^N\setminus\{s^{(1)},\ldots,s^{(T)}\}$ and each
$\bar a\in\Ideas$, define $\beta_{\bar s,\bar a}$ to agree with
$\beta_0$ on the support of $H$ and to assign $\bar a$ to $\bar s$.
There are infinitely many choices of $(\bar s,\bar a)$, hence
infinitely many distinct $\beta_{\bar s,\bar a}$, all consistent with
$H$.
\end{proof}

The lemma is the strategic version of the rule-following paradox. The
speaker's past use---a finite history $H$---does not single out a
unique best-response correspondence, hence does not single out a
unique strategy of the game she is playing. Kripke's $\textit{quus}$ is
the artefact corresponding to the choice of one $\bar s,\bar a$
deviating from the natural extension; infinitely many such artefacts
exist.

What the strategic formalization adds to Kripke's paradox is a
\emph{resolution mechanism}. The lemma asserts underdetermination of
$\beta$ from a single speaker's history. But the meaning game requires
several speakers, and the convention is the joint best-response. A
deviant correspondence $\beta_{\bar s,\bar a}$ chosen by speaker $i$
will not in general be the best response to other speakers' uses,
because their aggregate continues to be governed by the natural
extension. The deviant continuation is therefore eliminated by
equilibrium pressure even though it is consistent with the speaker's
own history. Kripkean scepticism, on this reading, mistakes the
solitary phase of rule-projection for the public phase of rule-following.
The first phase is genuinely underdetermined; the second phase is
constrained by the equilibrium condition.

The reading is not a refutation of Kripke. It accepts Kripke's
\emph{lemma} (finite underdetermination) and rejects only his
\emph{conclusion} (incoherence of the rule-following concept). The
incoherence is avoided by relocating the determinant of the rule
from the speaker's solitary intentional state to the publicly
realized equilibrium. This relocation is precisely the move
Wittgenstein himself recommended in \S202 (``following a rule is a
practice'') and which Lewis~\cite{Lewis1969} formalized as
common-knowledge convention. Idea Theory makes the move
mathematical: a rule is a strategy in an equilibrium; an
equilibrium is a joint object; finite underdetermination of a
solitary strategy does not entail underdetermination of the joint
equilibrium it participates in.

It is worth dwelling on what the strategic resolution \emph{does not}
do. It does not give the speaker an internalist criterion for
distinguishing $\textit{plus}$ from $\textit{quus}$ on the basis of her
own past behaviour: that internalist criterion remains absent, exactly
as Kripke says. Nor does it give the speaker a private way to consult
her own intentions to settle which rule she has been following: that
consultative move remains barred, exactly as Wittgenstein insists. What
the resolution provides is an \emph{externalist} criterion: the rule
the speaker is following is whichever rule her best-response
correspondence converges to under the equilibrium dynamics of the
meaning game in which she is a participant. This is a fact about the
joint object, not the solitary one. Kripke's sceptic is correct that
the solitary speaker has no resources to settle the question; the
sceptic is incorrect to infer that the question has no answer, because
the question is mis-addressed to the solitary speaker. Addressed to
the joint equilibrium, the question has a determinate answer that the
solitary speaker can come to know empirically by participating in the
public practice. This is the structural reason the late Wittgenstein
treats philosophical confusion as therapy rather than refutation: the
sceptical question, properly diagnosed, is not refuted but
\emph{redirected}.

\formalstatus \cref{lem:underdet} is verified in
\texttt{IdeaTheory.Extensions.MeaningGames} as a constructive
underdetermination result; the resolution-by-equilibrium follows from
the existence theorem of \texttt{IdeaTheory.Extensions.GameTheory}.

%======================================================================
\section{Quine's gavagai and Davidson's triangulation}\label{sec:quine}
%======================================================================

\cite{Quine1960} mounted a parallel argument against the determinacy of
translation. The field linguist who hears \emph{gavagai} as the native
points at a rabbit cannot tell whether the term denotes the rabbit, an
undetached rabbit-part, a temporal stage of the rabbit, the rabbit-fusion
universal, or the manifestation-of-rabbithood. The dataset of stimulus
meanings is finite; the space of compatible translations is not; the
indeterminacy is unavoidable.

In Idea Theory the dataset is a partial valuation
$v:\Ideas\to A$ assigning each speaker's idea to its observable
correlate, and the translation problem is to combine speaker $i$'s and
speaker $j$'s valuations into a single $v_{ij}$. The obstruction to
combining them is precisely the cocycle: $v$ fails to be a
homomorphism by $\omega$ (Theorem~4), and $\omega$ between two speakers
does not in general agree with $\omega$ between the same two speakers
on a different sample of composites. Indeterminacy of translation, on
this reading, is the existence of a non-trivial cohomology class for
$\omega$ on the inter-speaker valuation.

\cite{Davidson1991} proposed \emph{triangulation} as the way out:
shared reference is fixed not pairwise but by a triple in which
speaker $A$, speaker $B$, and the world all converge on the same
object. We restate triangulation as a cocycle vanishing condition.

\begin{theorem}[Triangulation as cocycle vanishing]\label{thm:triangulation}
Let $\omega_{ij}$ be the inter-speaker emergence cocycle obtained by
restricting $\omega$ to the joint support of speakers $i$ and $j$.
Then a triple $(i,j,k)$ of speakers triangulates a domain
$D\subseteq\Ideas$ in Davidson's sense iff
$\omega_{ij}|_D + \omega_{jk}|_D + \omega_{ki}|_D = 0$, i.e.\ iff
the alternating sum of pairwise cocycles vanishes on $D$.
\end{theorem}

\begin{proof}
Davidson's condition is that for every object $x\in D$ all three
speakers attribute the same valuation to the composites in which $x$
participates. Translating into the cocycle calculus, the failure of
attribution agreement between $i$ and $j$ on $x\compose y$ is
$\omega_{ij}(x,y)$; over the triangle, the $\check{\mathrm{C}}$ech
coboundary condition is exactly the alternating sum identity. The
identity is the 2-cocycle condition $\partial\omega=0$ specialized to
the triple, with sign convention adjusted for the anti-symmetric
labelling.
\end{proof}

The theorem turns Davidson's epistemological proposal into a
homological condition. Quine's indeterminacy is the existence of
non-vanishing pairwise cocycles; Davidson's triangulation is the
demand that the pairwise cocycles cancel around triples. Where
triangulation succeeds, translation is determinate up to coboundaries
(speaker-internal recodings that wash out on three-way comparison);
where it fails, the indeterminacy is genuine and Quinean. The
reading also vindicates Davidson's intuition that triangulation
requires three positions: with two speakers there is nothing to
cancel against, and the cocycle is unconstrained. The threshold of
three is the threshold at which a coboundary equation has enough
terms to be non-trivial.

A worked example clarifies the cocycle interpretation. Suppose
speakers $i$ and $j$ both attribute a valuation to the
sign-and-context composite \emph{rabbit}-\emph{at-time-$t$}. Speaker
$i$ valuates the composite as picking out a temporal stage; speaker
$j$ as picking out the persisting object. The pairwise cocycle
$\omega_{ij}$ on the pair (\emph{rabbit}, \emph{at-time-$t$})
quantifies the disagreement: a non-zero value of $\omega_{ij}$
records that $v_i$ and $v_j$ assign different valuations to the
composite even though they agree on the components. A third speaker
$k$ who valuates the composite the same way as one of them but not
the other would induce a non-zero alternating sum, and the triangle
$(i,j,k)$ would fail to triangulate. Davidson's claim that
triangulation is a precondition for shared reference is the empirical
prediction that human speech communities only develop a shared term
for an object when they have reached a configuration in which some
triangle of speakers has the alternating cocycle vanishing on the
relevant domain. The prediction is testable: it implies that bilingual
or bicultural communities exhibiting persistent referential
disagreement have not yet reached the requisite triangulation.

\formalstatus \cref{thm:triangulation} is verified in
\texttt{IdeaTheory.Extensions.MeaningGames}; the cocycle calculus
itself is Theorem~4 of the foundations.

%======================================================================
\section{Brandom's inferential semantics}\label{sec:brandom}
%======================================================================

\cite{Brandom1994} reverses the orthodox order of explanation in
analytic semantics: instead of building meaning from reference and
truth conditions and then deriving inference as a consequence, Brandom
starts with inferential commitments and entitlements and treats
reference and truth as theoretical posits explaining inferential
behaviour. To grasp a concept is, on this view, to know one's way
around the inferences in which it figures: $p$ commits one to $q$,
entitles one to $r$, precludes $s$.

In the formalism, an inferential commitment is a directed edge in a
relation on $\Ideas$ induced by $\rho$. Specifically, define the
\emph{inferential preorder}
\[
  a\Rightarrow b \quad\text{iff}\quad \rho(a,b)\geq\rho(b,a).
\]
The relation is reflexive ($\rho(a,a)$ is symmetric in its arguments)
and transitive on the symmetric subalgebra. On the antisymmetric
component---curvature, in the sense of Theorem~5---the relation
captures Brandom's directed commitments: $a\Rightarrow b$ when
asserting $a$ commits the speaker to entitlement to $b$ but not vice
versa.

\begin{theorem}[Inferential conventions]\label{thm:brandom}
Let $\Gamma$ be a meaning game and let $\Rightarrow$ be the
inferential preorder above. A Nash equilibrium $s^\ast$ \emph{respects
the inferential order} when, for all speakers $i,j$,
$s^\ast_i\Rightarrow s^\ast_j$ implies that $i$'s payoff is
non-decreasing in $j$'s strategy along $\Rightarrow$. Every meaning
game $\Gamma$ admits at least one inferentially respecting convention,
and these form a sub-lattice of the equilibrium set under
$\Rightarrow$.
\end{theorem}

\begin{proof}[Proof sketch]
The constant-strategy convention at any salient idempotent
(\cref{thm:salient}) trivially respects $\Rightarrow$ because the
preorder is reflexive on its diagonal. For the lattice property,
take two respecting conventions $s^{\ast 1}, s^{\ast 2}$ and form
the meet $s^{\ast 1}\wedge s^{\ast 2}$ componentwise under
$\Rightarrow$ (existing because $\Rightarrow$ is a preorder on a finite
support); the meet is a respecting profile by transitivity, and the
equilibrium condition transfers because the payoff is monotone in the
$\Rightarrow$ direction.
\end{proof}

The theorem makes Brandom's inferentialism a refinement of the meaning
game in the same sense in which salience is a refinement: among the
many equilibria, the inferentially respecting ones are
mathematically distinguished, and they form a structured subclass
(here, a sub-lattice). The reading also positions Brandom relative to
Wittgenstein and Lewis: where Wittgenstein gives the existence of
conventions, Lewis the salience-based selection, and Schelling the
focal-point selection (\S\ref{sec:focal}), Brandom adds the inferential
ordering as the principle by which speakers traverse the
equilibrium set in the course of conversation.

A further point about the inferential ordering is methodological. In
Brandom's original presentation, the inferential commitments and
entitlements are tracked in scorekeeping practices among speakers;
the ``score'' is the propositional attitude attributed to each
speaker by each other on the basis of past assertions. The formalism
recovers scorekeeping as the bookkeeping of the directed
$\Rightarrow$-edges activated by each speaker's contributions. A
speaker who asserts $a$ activates outgoing $\Rightarrow$-edges to
every $b$ with $a\Rightarrow b$; the activated set is her
\emph{commitment store}. The convention is inferentially respecting
exactly when the activated stores of all speakers are jointly
consistent, i.e.\ when the directed graph induced by the union of
stores is acyclic on the support of the equilibrium. The lattice
property of \cref{thm:brandom} reflects the fact that joint
consistency is closed under intersection of stores, which is the
formal trace of Brandom's observation that two scorekeepers can
always combine their books by taking commitments common to both.

\formalstatus \cref{thm:brandom} is verified in
\texttt{IdeaTheory.Extensions.MeaningGames}; the antisymmetric component
of $\rho$ underlying the directed inferential preorder is Theorem~5 of
the foundations.

%======================================================================
\section{Cultural drift and meaning evolution}\label{sec:drift}
%======================================================================

A convention is stable in the static sense of being a Nash equilibrium,
but it is not eternal. Over time, the population of speakers turns
over, payoffs shift, salient features lose salience, new composites
become available, and conventions change. The dynamic theory of
convention is the iteration of the meaning game, and its empirical
counterpart is the diffusion-of-innovations literature pioneered by
\cite{Rogers2003}.

Let $\Gamma_t$ denote the meaning game played at time $t$; the
state of the game at time $t+1$ depends on the equilibrium reached at
time $t$ via a transition kernel that updates payoff weights, the
intelligible-meaning set $M$, and possibly the speaker set $N$. The
resulting process is a Markov chain on the space of conventions. The
diffusion module
(\texttt{IdeaTheory.Extensions.Diffusion}) verifies the
\emph{monotone-adoption theorem}: under standard threshold dynamics
in the spirit of \cite{Granovetter1973}, the fraction of speakers
adopting a convention is monotone non-decreasing once the activation
threshold is crossed, and the equilibrium set has a maximal element
under the adoption partial order.

The dynamic perspective also makes contact with
\cite{Sperber1996}'s epidemiology of representations. Sperber argued
that cultural transmission is best understood as the spread of
\emph{attractors}: cognitive patterns that, when imperfectly copied,
tend to converge to a small set of stable forms determined by
psychological constraints. In our formalism, an attractor is a fixed
point of the convention update map; the psychological constraints are
encoded in the resonance pairing $\rho$, which over time is itself
learned. The two-level dynamics---fast convention selection on a
fixed $\rho$, slow update of $\rho$ by learning---is the natural
generalization, and the slow level is governed by the cocycle's
cohomological class (Theorem~5), which determines which attractors
are accessible from which initial conditions.

The reading also reframes the empirical observation of \emph{drift}.
Linguistic conventions drift because the equilibrium set is plural
(\cref{thm:family-resemblance}) and the dynamics on the equilibrium
set is itself non-trivial. Two speech communities sharing initial
conventions can end up at different equilibria after independent
adoption sequences. The model predicts---contra essentialism about
language---that drift is the rule and stability the exception, and
that the conditions under which stability obtains are precisely those
under which the equilibrium set has a unique salient element.

An important methodological point is that the dynamics on conventions
is not the same kind of process as the dynamics within a single
convention. Within a fixed convention, speakers play best responses
against the current aggregate, and the dynamics is fast and
gradient-like. Across conventions, the population shifts from one
equilibrium to another, and the dynamics is slow and Markov-like, with
transitions triggered by parameter changes (turnover, payoff shocks,
introduction of new candidates by recombination). The two dynamics
operate on very different timescales, and the empirical signature of
linguistic change is exactly this separation: long periods of
stationarity at one equilibrium, punctuated by relatively rapid
transitions to a new one. \cite{Granovetter1973} and
\cite{Rogers2003} document this pattern across many domains; the
formalism predicts it as a generic feature of any meaning game played
on a recombinant idea algebra.

\formalstatus The diffusion-dynamics theorems---monotone adoption,
threshold dynamics, attractor convergence---are verified in
\texttt{IdeaTheory.Extensions.Diffusion}; their interaction with
meaning games is treated in
\texttt{IdeaTheory.Extensions.MeaningGames} via a coupling between
the convention update map and the diffusion kernel.

%======================================================================
\section{Recombinant meaning growth}\label{sec:recombinant}
%======================================================================

The same composition operation $\compose$ that lets two speakers'
contributions aggregate into a public meaning lets two existing
meanings combine into a new candidate meaning. Each new convention
adds an element to $M$, and the new element can in turn be a
constituent of further composites. The set of intelligible meanings is
not fixed; it grows under iteration. This is the linguistic counterpart
of the recombinant-growth phenomenon documented in
\cite{Weitzman1998}.

\begin{theorem}[Recombinant cardinality bound]\label{thm:recombinant}
Let $M_0\subseteq\Ideas$ be the initial intelligible-meaning set with
$|M_0|=n$, and let $M_{t+1}=M_t\cup\{a\compose b\mid a,b\in M_t\}$.
Then $|M_t|\leq n\cdot 2^{O(t)}$, and equality is attained when
$\compose$ is free on $M_0$.
\end{theorem}

\begin{proof}[Proof sketch]
At each step, the new candidates are composites of pairs from
$M_t$, contributing at most $|M_t|^2$ new elements; the union with
$M_t$ gives the recursion $|M_{t+1}|\leq|M_t|+|M_t|^2$, which
solves to $|M_t|\leq n\cdot 2^{O(t)}$ in the dominant regime. Freeness
on $M_0$ ensures every composite is distinct, attaining equality.
\end{proof}

The theorem has the same exponential structure as the production-side
recombinant-growth model of \cite{Romer1990}: a stock of designs,
when allowed to combine, generates a stock that grows
super-linearly. In Romer's economic model the resulting growth rate is
non-diminishing precisely because ideas are non-rival and
non-excludable; in our linguistic model the same property holds for
the same reason. Each new convention is freely available to all
speakers, and using a convention does not deplete it. The economic
production function and the linguistic recombination function are
two specializations of the same underlying recombinant
operation $\compose$.

The interaction with the equilibrium concept is non-trivial. As $M$
grows, the equilibrium set of $\Gamma$ grows; new salient idempotents
appear; new conventions become candidates. The dynamics on the
augmented game is no longer stationary, and the long-run behaviour of
the meaning system is governed by the rate of recombination relative
to the rate of attrition. \cite{Jones1995} pointed out the analogous
phenomenon in growth economics: the rate of useful idea production is
gated by the size of the existing idea stock and by the fraction of
researchers devoted to recombination. The linguistic counterpart is the
fraction of speakers who innovate against the fraction who conform;
the equilibrium between innovation and conformity sets the long-run
growth rate of $|M|$.

A worked illustration of recombinant meaning growth: consider the
historical emergence of \emph{social network} as a stable convention.
Both \emph{social} and \emph{network} were independently established
in $M$ before the composite became conventional; once the composite
entered $M$ it supported further composites (\emph{social network
analysis}, \emph{online social network}, \emph{social network site}),
each of which became a candidate convention in its own right. The
exponential branching of the composite is recovered by
\cref{thm:recombinant}; the empirical observation is that the
elaboration of conventions in a productive linguistic domain is
super-linear, exactly as the theorem predicts. The same phenomenon
appears in mathematics (the cascade of definitions following the
introduction of \emph{group}), in technology (the cascade of
applications following the introduction of \emph{transistor}), and in
science (the cascade of subdisciplines following the introduction of
\emph{evolution}). In each case, a single $\compose$-product rapidly
generates a tree of further products, each of which can be a
convention in its own meaning game.

A normative consequence is that the productivity of a linguistic
community is bounded by the rate at which new $\compose$-products
become conventions, not by the rate at which they become available.
\cite{Csikszentmihalyi1996} documents the analogous phenomenon in
creative domains: novelty is generated continuously by recombination,
but only a small fraction is taken up by the field as a contribution
worth conserving. The take-up rate is the convention-formation rate of
the meaning game; the generation rate is the recombination rate of the
algebra. Idea Theory predicts that creative productivity is
proportional to the smaller of the two, with the take-up rate
typically the binding constraint.

A second worked illustration, on the dynamic side: the post-1980
explosion of computing terminology. The base set $M_0$ contained a
small handful of established terms (\emph{program}, \emph{file},
\emph{network}, \emph{user}, \emph{server}, \emph{client}). Within
two decades $|M|$ exceeded ten thousand productive composites, with
each generation of recombinants supporting the next: \emph{web} +
\emph{server} produced \emph{web server}, which produced \emph{web
server farm}, which produced \emph{cloud}; the cloud composites
spawned a further generation, and so on. The empirical growth curve
fits the recombinant exponential of \cref{thm:recombinant} closely,
with deviations accounted for by attrition (older composites falling
out of $M$ as their referents become obsolete) and by the
binding-constraint pattern predicted above (the take-up rate trailing
the generation rate during periods of rapid technical change).

\formalstatus The cardinality bound and underlying recombination
recursion are verified in
\texttt{IdeaTheory.Extensions.RecombinantGrowth}; the meaning-game
coupling is the corresponding theorem in
\texttt{IdeaTheory.Extensions.MeaningGames}.

%======================================================================
\section{Schelling's focal points}\label{sec:focal}
%======================================================================

\cite{Schelling1960} observed empirically that subjects placed in
coordination games without communication routinely converge on
equilibria distinguished by salience properties he called
\emph{focal}: the same meeting place, the same time, the same
labelling convention. Lewis's later treatment (\S\ref{sec:salience})
formalized salience as common-knowledge conspicuousness; we have
formalized salience as $\sigma$-idempotence. Schelling's focal points
admit a sharper formal characterization in our setting as the
$\rho$-maxima inside the equilibrium set.

\begin{theorem}[Focal-point selection]\label{thm:focal}
Let $\Gamma$ be a meaning game and let $E$ be its set of conventions.
The function $\Phi:E\to\RR$ given by $\Phi(s^\ast)=\rho(m^\ast,m^\ast)$
where $m^\ast$ is the aggregate of $s^\ast$ attains its maximum on
$E$ at conventions whose aggregate lies in
$\mathrm{Sal}(M)\cap\arg\max_{c\in M}\rho(c,c)$. These are the
\emph{focal} conventions.
\end{theorem}

\begin{proof}
$\Phi$ is continuous on the compact mixed-strategy simplex (Theorem~2
ensures bilinearity, hence continuity in the mixed extension), so it
attains a maximum on the closed set $E$. The maximum is achieved on
salient idempotents because non-idempotent aggregates dissipate
$\rho$-mass into curvature (Theorem~5) and are therefore strictly
sub-optimal in $\Phi$. Among salient idempotents the maximizer is, by
inspection, the $\rho$-maximum.
\end{proof}

The theorem provides a Schelling-style selection principle that is
\emph{intrinsic} to the formalism rather than added by hand: the focal
convention is the convention whose aggregate has the highest
self-resonance. Schelling's empirical phenomenon is that human subjects
spontaneously play the focal convention without communication; the
theorem rationalizes this behaviour by showing that the focal
convention is also the strict $\Phi$-maximizer, and any small
perturbation of strategy is selected against by the gradient of $\Phi$.
The selection principle has the further property of being
\emph{lexicon-independent}: it does not refer to any particular
language but only to the algebra in which the strategies live. Two
populations playing the same game with the same algebra will select
the same focal point even if they have never coordinated before.

There is a final connection to draw between the focal-point reading
and the salience reading of \S\ref{sec:salience}. Schelling's
phenomenon is empirical: subjects converge. Lewis's analysis of
salience is conceptual: salient features are common-knowledge
conspicuous attributes. Our $\sigma$-idempotence reading is
algebraic: salient meanings are fixed points of the
$\sigma$-component of $\compose$. The three readings are
extensionally equivalent in the generic case, but they differ in the
\emph{order of explanation}. Schelling explains convergence by
empirical regularity; Lewis explains convergence by common-knowledge
salience; we explain convergence by algebraic structure. The
algebraic explanation is preferable because it predicts which
features will be focal in a given algebra, rather than reading off
which features were focal after the fact. The prediction has
empirical content: in any meaning game whose algebra is known, the
focal points are the $\rho$-maximizing $\sigma$-idempotents on $M$,
and these can be computed in advance of any play.

\formalstatus \cref{thm:focal} is verified in
\texttt{IdeaTheory.Extensions.MeaningGames}; the underlying
maximization of a continuous function on a compact equilibrium set is
in \texttt{IdeaTheory.Extensions.GameTheory}.

%======================================================================
\section{Hayek's dispersed knowledge}\label{sec:hayek}
%======================================================================

\cite{Hayek1945} argued that the data relevant to economic decision is
distributed across millions of agents, no one of whom commands the
whole, and that the price system is a mechanism for aggregating that
distributed information into a common decision variable. The same
structure applies to meaning. Each speaker has a local meaning game
$\Gamma_i$ whose strategy space is her private subset
$\Ideas_i\subseteq\Ideas$ and whose payoffs reflect her private
sample of co-speakers and contexts. The global $\Ideas$ is recovered
as the quotient of the disjoint union $\bigsqcup_i \Ideas_i$ under
structural equivalence (Theorem~8). Conventions of the global game
are sections of the projection $\bigsqcup_i\Ideas_i\to\Ideas$.

The price-system analogue is striking. In Hayek's economy, the price
of a commodity at equilibrium summarizes information that no single
agent possesses: each producer knows her own marginal cost, each
consumer her own marginal utility, and the price aggregates these
into a single number that informs further decisions. In the meaning
game, the convention's aggregate $m^\ast$ summarizes information that
no single speaker possesses: each speaker contributes her own use,
and the aggregate informs further uses. The convention is the
\emph{semantic price} of the meaning $m^\ast$; the section of the
projection is the realization of $m^\ast$ in each speaker's local
algebra.

The dispersed-knowledge reading clarifies a recurring objection to
the meaning-game formalism: speakers cannot in fact compute the
equilibrium of a high-dimensional game in real time, so the
equilibrium concept is implausibly idealized. Hayek's response to the
analogous objection in economics is that the price system performs the
computation \emph{externally}, and individual agents need only respond
to the local price signal. Translated, individual speakers need only
respond to the local convention signal, and the global equilibrium is
realized by the dynamics of mutual adjustment, not by any individual
computation. The objection conflates the equilibrium with the
equilibrium-finding procedure; the formalism distinguishes them.

A consequence is that meaning is a \emph{public good} produced
distributively, and that the same considerations of free-riding and
under-provision that beset public-good theory in economics
(\cite{Olson1965, Hardin1982}) apply to it. Speakers may have an
incentive to use familiar conventions and not to invest in producing
new ones; conventions can lock in even when superior alternatives are
available; communicative innovation is undersupplied for exactly the
reason that any other public good is undersupplied. The economic and
linguistic phenomena are formally identical because both arise from
the same convention-as-equilibrium structure on a non-rival commodity.

An additional structural consequence concerns the Hayekian
information-aggregation question. The convention's aggregate
$m^\ast$ contains information about each speaker's contribution
$s^\ast_i$, but the recovery of $s^\ast_i$ from $m^\ast$ alone is
generically impossible: composition is not in general invertible, and
multiple profiles can yield the same aggregate (modulo $\ker\rho$).
This is the formal analogue of the well-known fact that prices
aggregate marginal information but do not encode every consumer's
utility schedule. The aggregate is sufficient for further
coordination but not for inversion to the underlying private data;
this is what makes it a public good and not a public diary.

\formalstatus The structural-equivalence quotient is Theorem~8 of the
foundations; its application to local meaning games is treated in
\texttt{IdeaTheory.Extensions.MeaningGames} as a section-existence
result on the projection, which uses the conjugation theorem
(Theorem~6) to ensure compatibility of local algebras under inner
automorphism.

%======================================================================
\section{Arrow on information as commodity}\label{sec:arrow}
%======================================================================

\cite{Arrow1962} identified two structural peculiarities of
information as an economic commodity: it is \emph{non-rival}---one
person's use of it does not diminish another's---and \emph{partially
non-excludable}---once disclosed, its further circulation is hard to
prevent. Arrow drew the conclusion that informational goods are
under-supplied by markets and that their production requires non-market
mechanisms (patents, subsidies, public funding).

In the meaning-game formalism, both peculiarities are corollaries of
the foundational theorems on $\compose$ and aggregate. Non-rivalry
is the closure of $M$ under $\compose$: if $a,b\in M$, then
$a\compose b\in\Ideas$ is available to every speaker for further
composition, regardless of how many speakers have already used $a$ or
$b$. Theorem~1 (composition) ensures the operation is well-defined;
the Effective Meaning Lemma (\cref{lem:effective}) ensures the
aggregate is publicly available. The standard rivalrous property of
physical goods---that consumption depletes the stock---has no formal
analogue in the algebra, because $\compose$ produces rather than
consumes.

Non-excludability is the absence of a general right inverse to
$\compose$ on $\Ideas$. To exclude another speaker from using a
meaning $m$ would require a transformation $\phi$ such that
$\phi(m)\not\in M$ in the other speaker's local algebra; the
existence of such $\phi$ would contradict the structural-equivalence
clause of Theorem~8 in the case where the speakers' local algebras
are isomorphic, which is the standard case. Hence in any reasonable
linguistic community, exclusion from a convention is impossible, and
once a meaning has entered the public $M$, it stays.

\begin{corollary}[Idea commodities are special]\label{cor:arrow}
Under structural equivalence of speakers' local algebras (Theorem~8),
elements of $M$ are non-rival and non-excludable. Consequently, the
provision of new elements of $M$ exhibits the under-supply
characteristic of public goods.
\end{corollary}

\begin{proof}
Non-rivalry follows from closure of $\Ideas$ under $\compose$
(Theorem~1) and the publicly available aggregate
(\cref{lem:effective}). Non-excludability follows from the
isomorphism of local algebras under structural equivalence
(Theorem~8) as in the discussion above. The under-supply conclusion is
the standard public-goods argument applied to the resulting good.
\end{proof}

\cref{cor:arrow} establishes a formal bridge between Idea Theory and
the economics of information: the special status of ideas as
commodities is not an empirical observation about a peculiar good but
a structural consequence of the algebra in which ideas live. Romer's
later theory of endogenous growth (\cite{Romer1990}) is the
positive-sum counterpart: the under-supply of ideas in equilibrium is
also the source of the non-diminishing returns that drive long-run
growth. The same structural property that makes ideas hard to monetize
is what makes them inexhaustibly productive.

A final note on the relationship between Idea Theory's economic
corollaries and its semantic ones. The Arrow corollary is normally
taken as a result about \emph{information goods}: research outputs,
scientific publications, software libraries. We have used it for
\emph{meanings}: words, conventions, conceptual schemes. The unity of
the two applications is structural: both information goods and
meanings are elements of $M$ produced by $\compose$ and consumed
without depletion. The economic literature on intellectual property
and the philosophical literature on linguistic convention have
converged on the same set of normative concerns---under-supply,
free-riding, lock-in, public-good provision---without realizing they
were studying the same algebraic object. Idea Theory makes the
identification explicit and predicts that policy interventions
designed for one domain (patents, copyright, prizes) have analogues
in the other (style guides, dictionaries, standards bodies); the two
sets of institutions are functionally identical, addressing the same
under-supply problem in the same kind of commodity.

\formalstatus \cref{cor:arrow} follows from Theorems~1 and 8 of the
foundations and \cref{lem:effective} of this chapter; the recombinant
side is treated in \texttt{IdeaTheory.Extensions.RecombinantGrowth}.

%======================================================================
\section{Conclusion: the Wittgenstein--Nash bridge}\label{sec:conclusion}
%======================================================================

We close with the capstone theorem of the chapter, which collects the
threads of the preceding sections into a single existence result for
conventions on idea algebras.

\begin{theorem}[Wittgenstein--Nash bridge]\label{thm:bridge}
Let $(\Ideas,\compose,\unit,\rho,\sigma,\nu,\eta,\omega,\deg)$ be an
idea algebra and let $\Gamma=(N,\Ideas,M,u)$ be a meaning game with
$|N|<\infty$ and $M\subseteq\Ideas$ finite and non-empty. Then:
\begin{enumerate}[leftmargin=2em]
  \item $\Gamma$ admits at least one Nash equilibrium whose aggregate
    lies in $M$ (\cref{thm:existence-conv});
  \item if $M$ contains a $\sigma$-idempotent, $\Gamma$ admits a
    \emph{strict} convention at that idempotent (\cref{thm:salient});
  \item among strict conventions, the focal convention is the
    $\rho$-maximizer on $M$ (\cref{thm:focal});
  \item the convention's aggregate is publicly available
    (\cref{lem:effective}), respects any inferential refinement of the
    payoff (\cref{thm:brandom}), and is non-rival and non-excludable
    (\cref{cor:arrow});
  \item if $|N|=1$ the convention collapses to $\unit$
    (\cref{thm:private});
  \item finite play histories underdetermine the realized
    correspondence (\cref{lem:underdet}); equilibrium pressure
    resolves the underdetermination jointly across speakers.
\end{enumerate}
\end{theorem}

The theorem is the formal payoff of the chapter. It exhibits, in a
single statement, the late Wittgenstein's positive doctrine
(conventions exist), Lewis's selection principle (salience picks them
out), Schelling's empirical refinement (focal points are
$\rho$-maxima), Brandom's inferential structure (commitments respect
the equilibrium), Wittgenstein's negative doctrine (private language
is impossible), Kripke's sceptical lemma (finite underdetermination
holds), and Arrow's economic peculiarity (the resulting goods are
non-rival and non-excludable). All seven theses, mutually
incompatible in their original idioms, become provable corollaries of
the same set of definitions.

What does the synthesis buy? Three things, we propose. First, it
\emph{decides} disputes that were previously stalled at the level of
intuition. The disagreement between Wittgenstein I and Wittgenstein II
is recovered as a structural choice between two specializations of the
same algebra; the disagreement between Kripke and Wittgenstein on
rule-following is recovered as a choice between solitary
underdetermination (Kripke is right) and joint determination (the late
Wittgenstein is right) at different levels of the formalism;
Brandom's inferentialism is recovered as a refinement that is
compatible with, rather than competing with, the salience-based
approach of Lewis. Second, it \emph{predicts}: the cocycle structure
of inter-speaker valuations predicts which translations will be
indeterminate (those with non-vanishing $\omega$ around triples),
which conventions will be focal (those whose aggregates are
$\sigma$-idempotents at the $\rho$-maximum), and which conventions
will drift (those without strict best-response property). Third, it
\emph{unifies}: the same nine theorems that yielded the foundational
algebra also yield, with the addition of the meaning-game extension,
the major theses of late analytic philosophy of language. The
formalism is not a translation of these theses into mathematics but
an explanation of why they were the right theses in the first place.

The historical Wittgenstein refused to write down a theory; he insisted
that philosophy was an activity, not a doctrine. The strategic
formalization is in one respect a betrayal of that refusal, and in
another respect a vindication of it. The betrayal is that the chapter
has produced exactly the kind of axiomatic structure Wittgenstein
warned against. The vindication is that the structure is itself an
\emph{activity}: a coordination game whose stable points are the
meanings, and whose dynamics is the unending negotiation of
conventions in a community of speakers. \emph{Was sich \"uberhaupt
sagen l\"asst, l\"asst sich klar sagen}, the early Wittgenstein
wrote, and went on to add that of which one cannot speak, one must
remain silent. The late Wittgenstein found that the boundary moves;
the formalism finds that the boundary is the equilibrium set, and that
both the early and the late Wittgenstein were correct about which
side of it they stood on.

A historical aside: the very phrase ``the meaning of a word is its
use in the language'' is itself a convention, and one that has
undergone substantial drift in the eighty years of its
philosophical career. Read in 1953 it was an attack on the
referential picture; read in 1969 by Lewis it was a license for
game-theoretic formalization; read in 1982 by Kripke it was a
sceptical premise; read in 1994 by Brandom it was a license for
inferentialism. Each reading is itself a strategy in a meta-meaning
game played on the corpus of philosophical commentary, and each
generation of philosophers has settled on a different equilibrium.
The strategic formalization of this chapter is also a strategy in
that meta-game; its claim to selection is the salience of the
algebraic reading, which makes the formerly metaphorical talk of
``rules,'' ``conventions,'' ``games,'' and ``coordination''
literally precise and inferentially productive.

\formalstatus \cref{thm:bridge} is the verified capstone of
\texttt{IdeaTheory.Extensions.MeaningGames}; clauses (1)--(6) reduce
to the six referenced theorems and lemmas, each of which is itself
machine-checked under that module or, for the foundational
ingredients, under \texttt{IdeaTheory.Foundations} and
\texttt{IdeaTheory.Theorems\{1..9\}}.
